% =============================================================================
% Shared preamble of the solutions to the real-exercise list for Lectures 1--5.
% Input by main.tex (all lectures) and lecture-0N.tex (one lecture). The driver defines
% \soltema (running head) before \input-ing this file. Layout copied from
% Listas/lista-reais-lectures-01-04-solutions/main.tex.
% Compile from this folder: pdflatex <driver>.tex (twice); pdffonts <driver>.pdf (all Type 1).
% =============================================================================
\usepackage[utf8]{inputenc}
\usepackage[T1]{fontenc}
\usepackage{lmodern}   % Type 1 vector fonts (replaces PK bitmaps)
\usepackage{cmap}      % ToUnicode CMap -> copyable text and accents
\usepackage[english]{babel}
\usepackage{amsmath,amssymb,amsthm,mathtools}
\usepackage[a4paper,margin=2.4cm]{geometry}
\usepackage{microtype}
\usepackage{booktabs,tabularx,array}
\usepackage{enumitem}
\usepackage{xcolor}
\usepackage{tikz}
\usetikzlibrary{arrows.meta,calc}
\usepackage[most]{tcolorbox}
\usepackage{fancyhdr}
\usepackage[hidelinks]{hyperref}

\setlist{topsep=3pt,itemsep=2pt}
\setlength{\parindent}{0pt}
\setlength{\parskip}{5pt}
\renewcommand{\arraystretch}{1.2}

\pagestyle{fancy}\fancyhf{}
\lhead{\footnotesize Economics of Organizations and Incentives --- MPE Insper}
\rhead{\footnotesize \soltema}
\cfoot{\footnotesize\thepage}
\renewcommand{\headrulewidth}{0.3pt}
\setlength{\headheight}{14pt}

\definecolor{hdr}{RGB}{30,80,150}
\definecolor{ans}{RGB}{20,110,70}

% ----------------------------------------------------------------------------- numbering
% Item files carry no numbers. A driver calls \solitem{<combined n>}{<slug>}: it steps the
% per-lecture counter ItemNo, sets \CombNo and inputs sol-<slug>.tex (skipped with a log
% message if the file does not exist yet). \sollecture{<title>} opens a lecture and resets
% ItemNo, so items are numbered 1--k within each lecture in every driver.
\newcounter{ItemNo}
\newcommand{\CombNo}{0}
\newcommand{\solitem}[2]{\stepcounter{ItemNo}\renewcommand{\CombNo}{#1}%
  \IfFileExists{sol-#2.tex}{\input{sol-#2}}{\typeout{solitem: sol-#2.tex missing, skipped}}}
% The page break sits here, so a lecture's first item follows its heading on the same page.
\newif\ifafterlecture
\newcommand{\sollecture}[1]{\clearpage\phantomsection\part*{#1}\addcontentsline{toc}{part}{#1}%
  \setcounter{ItemNo}{0}\afterlecturetrue}

% \begin{exsol}{<source and exercise>}{<printed page>}{<level>}{<difficulty 1-5>} ... \end{exsol}
% One exercise of the list; the heading shows its number in the lecture and in the combined list.
\newenvironment{exsol}[4]{%
  \ifafterlecture\global\afterlecturefalse\else\clearpage\fi\phantomsection
  \addcontentsline{toc}{section}{Item \theItemNo\ (\CombNo): #1 (p.~#2)}%
  \section*{\large Item \theItemNo\ {\normalsize\mdseries(item \CombNo\ of the combined list)}:
    #1\hfill{\normalsize\mdseries #3 $\cdot$ difficulty #4/5}}%
  \textbf{Source:} #1, page~#2. The statement is reproduced verbatim in the list (item
  \CombNo); it is not repeated here.\par}{\par}
% Fixed sub-heads, in this order inside each exsol.
\newcommand{\asked}{\paragraph{(a) What is asked.}}
\newcommand{\method}{\paragraph{(b) Method.}}
\newcommand{\solution}{\paragraph{(c) Solution.}}
\newcommand{\interp}{\paragraph{(d) Interpretation.}}
\newcommand{\verify}{\paragraph{(e) Verification.}}
\newcommand{\mistakes}{\paragraph{(f) Common mistakes.}}
% Boxed final answer: \answer{...}
\newcommand{\answer}[1]{\begin{tcolorbox}[colback=ans!5,colframe=ans,boxrule=0.6pt,
  left=5pt,right=5pt,top=3pt,bottom=3pt]\textbf{Answer.} #1\end{tcolorbox}}
% Best-response mark in a payoff matrix: \br{8} = payoff 8 is a best reply of its player.
\newcommand{\br}[1]{\underline{#1}}
\newcommand{\MG}{Mu\~noz-Garcia}
\newcommand{\Sm}{S_{-i}}
\newcommand{\sm}{s_{-i}}
